\documentclass[12pt,a4paper]{report}

\usepackage[utf8]{inputenc}
\usepackage[french]{babel}
\usepackage[autolanguage]{numprint}
\usepackage[T1]{fontenc}
\usepackage{amsfonts,amsmath,amssymb}
\usepackage{bigints}

\allowdisplaybreaks

\DeclareMathOperator{\arcsinh}{arcsinh}
\DeclareMathOperator{\arccosh}{arccosh}
\DeclareMathOperator{\arctanh}{arctanh}
\DeclareMathOperator{\cotan}{cotan}

\usepackage[pdftex]{pict2e}
\usepackage[dvipsnames]{xcolor}
\usepackage{graphicx}

\graphicspath{
    {chapters/graphics/}
	{chapters/exercices/graphics/}
}

\setcounter{tocdepth}{3}

\setlength{\paperwidth}{21cm}
\setlength{\paperheight}{29.7cm}
\setlength{\evensidemargin}{-1.0cm}
\setlength{\oddsidemargin}{-1.5cm}
\setlength{\topmargin}{-2cm}
\setlength{\headsep}{0.15cm}
\setlength{\headheight}{0.7cm}
\setlength{\textheight}{26cm}
\setlength{\textwidth}{19cm}

\newcommand{\be}{\begin{equation}}
\newcommand{\ee}{\end{equation}}
\newcommand{\benn}{\begin{equation*}}
\newcommand{\eenn}{\end{equation*}}
\newcommand{\bea}{\begin{eqnarray}}
\newcommand{\eea}{\end{eqnarray}}

\newcommand{\trint}{\mathop{\int\!\!\!\int\!\!\!\int}\limits}
\newcommand{\dint}{\mathop{\int\!\!\!\int}\limits}

\title{\'Etapes de calcul, d\'emonstrations et exercices du livre << Introduction au calcul tensoriel >> de Claude Senay et Bernard Silvestre-Brac}
\author{S\'ebastien Majerowicz}
\date{2026 -- 202X}

\begin{document}

\maketitle

\begin{abstract}
	Ce document est un condens\'{e} du livre en titre mais aussi une aide pour les d\'{e}monstration et la correction des exercices fournis.

	Une grandeur physique poss\`{e}de une existence intrins\`{e}que, ind\'{e}pendante de tout syst\`{e}me de r\'{e}f\'{e}rence.

	Le calcul tensoriel est la formalisme math\'{e}matique qui permet d'exprimer les lois de la physique sous une forme ind\'{e}pendante du syst\`{e}me de r\'{e}f\'{e}rence choisi.
\end{abstract}

\tableofcontents
\listoffigures

\clearpage
\pagenumbering{arabic}
\setcounter{page}{1}

\chapter{Rappels et conventions d'\'{e}criture}

\section{Conventions d'Einstein}

Une suite \emph{contravariante} est l'ensemble de $n$ quantit\'{e}s \'{e}tiquet\'{e}es par la m\^{e}me lettre et qui diff\`{e}rent par un indice plac\'{e} en position sup\'{e}rieure : $\{a^{1},a^{2},\dots,a^{n}\}$. Une suite \emph{covariante} est construire de mani\`{e}re identique \`{a} l'exception du fait que l'indice est en position inf\'{e}rieure : $\{b_{1},b_{2},\dots,b_{n}\}$.


\end{document}
